\section{Дикость построенных выпрямителей}\label{sec:group-results}

В этом разделе $K=\Q[a]/(a^2-a+1)$, $b=1-a$, $R=K[x]$.
Используются равенства $ab=1$ и $b^2=-a$.
Ручная группа $\Tame_3(K)$ порождена аффинными и элементарными
полиномиальными автоморфизмами. Координата называется ручной, если она
входит в ручной автоморфизм, и дикой в противном случае.
Относительная ручность над $R$ означает ручность в двух переменных над
кольцом $R$; это условие первоначально сильнее ручности над полем
коэффициентов $K$. Приведённые ниже выводы из опубликованных теорем
являются рабочими доказательствами этой исследовательской ветки и
предлагаются для дальнейшей проверки.

\begin{proposition}[Выпрямитель Шастри]\label{prop:trefoil-wild}
Координата $h_T$ и построенный выпрямитель $\Phi_T$ являются дикими
над $K$.
\end{proposition}
\begin{proof}
В относительных координатах $(u,v)=(axy+z,-y)$ базовая координата равна
$H=mv+P_x(u)$, где $m=bx^5$, а коэффициент при $u^2$ равен $bx$.
Критерий Эдо--Венеро для координат длины один
\cite[предложение~2]{EV2001} требует для ручности делимости всех
нелинейных коэффициентов на $m$. Но $bx^5\nmid bx$ в $R$.
Замена $(y,z)\mapsto(u,v)$ ручная над $R$, поэтому $H$ дикая над $R$.

По теореме Умирбаева--Ю дикая $K[x]$-координата остаётся дикой
координатой $K[x,y,z]$ в характеристике нуль
\cite[главная теорема]{UY2004}.
Предварительная замена $M=(x,y+2,z-5x)$ ручная, поэтому $h_T=H\circ M$
тоже дикая. Ручный автоморфизм не может иметь дикую компоненту;
следовательно, $\Phi_T$ дикий.
\end{proof}

\subsection{Внутренняя координата рациональной длины два}

Для следующей леммы $Y,Z$ обозначают две переменные над $R$.
Запись
\[
 F=p_2Y+Q_2(p_1Z+Q_1(Y))
\]
называем приведённой в смысле Эдо, если $p_1$ монический,
$Q_1\in R[Y]$, $Q_1(0)=Q_2(0)=0$,
$\gcd(p_1,\text{коэффициенты }Q_1)=1$;
здесь $p_1,p_2\ne0$, $p_2\in K(x)$, $Q_2\in K(x)[Y]$.
Мы используем случай $\deg Q_1\ge2$, $\deg Q_2\ge2$.

\begin{lemma}[Следствие классификации Эдо]\label{lem:inner-wild}
Пусть $F\in R[Y,Z]$ является координатой и имеет указанную
приведённую запись. Если $F$ ручная над $R$, то
$p_1Z+Q_1(Y)$ также ручная над $R$.
\end{lemma}
\begin{proof}
Здесь композиции $\sigma\tau$ понимаются как композиции
$R$-алгебраических автоморфизмов: $(\sigma\tau)(f)=\sigma(\tau(f))$.
Теорема~8(1) Эдо даёт $F=\sigma\tau(Y)$, где обе компоненты каждого
из $\sigma,\tau$ имеют длину один; такие автоморфизмы ручные.
Как в доказательстве теоремы~8(3), нормируем
$\sigma(Y)=p_3Z+Q_3(Y)$, где $p_3$ монический и $Q_3(0)=0$.
Нормировка аффинная и сохраняет ручность.
Формула компаньона Рассела--Сатайе имеет вид
\[
 \sigma(Z)=\frac{\varepsilon}{p_3}
 \bigl(V(p_3Z+Q_3(Y))-Y\bigr),\qquad
 \varepsilon\in R^\times,\quad V\in R[Y].
\]
Если $\tau(Y)=p_4Z+Q_4(Y)$, то
\[
 F=-\frac{\varepsilon p_4}{p_3}Y+
 \left(\frac{\varepsilon p_4}{p_3}V+Q_4\right)
 (p_3Z+Q_3(Y)).
\]
После удаления постоянного слагаемого это приведённая запись того же
полинома. Её единственность, установленная в предложении~2 Эдо,
даёт $p_3=p_1$, $Q_3=Q_1$. Поэтому внутренний полином является
компонентой ручной $\sigma$; см.\ также ту же нормировку в
доказательстве теоремы~8(3) \cite{Edo2013}.
\end{proof}

\begin{proposition}[Гармоническая модель и все построенные $C_N$]
\label{prop:family-wild}
Для каждого $N\ge1$ координата $H_N$ и построенный выпрямитель
$\Phi_N$ дикие над $K$. В частности, это верно для гармонической
модели $C_1=\gamma_H$.
\end{proposition}
\begin{proof}
Положим
\[
 Q_x(U)=bxU^2+(-1+ax^2)U+g(x).
\]
Из $P_x(U)=Q_x(U)-ax^4U$ и $a^2+b=0$ получаем тождество
\[
 H_N=H(x,p_N+x^Nz,y)
     =-ax^4y+Q_x(ax^{N+1}z+axp_N+y).
\]
Введём
\[
 \begin{aligned}
 G_N&=x^{N+1}z+by+x\bigl(p_N(x,y)-p_N(x,0)\bigr),\\
 c_N&=axp_N(x,0).
 \end{aligned}
\]
Тогда
\begin{equation}\label{eq:family-reduced}
 H_N-Q_x(c_N)
 =-ax^4y+\bigl[Q_x(aG_N+c_N)-Q_x(c_N)\bigr].
\end{equation}
Это приведённая запись из леммы~\ref{lem:inner-wild}:
$p_1=x^{N+1}$, $p_2=-ax^4$,
$Q_1(y)=by+x(p_N(x,y)-p_N(x,0))$.
Линейный коэффициент $Q_1$ равен $b$ по модулю $x$, поэтому
условие взаимной простоты выполнено. Внешний полином имеет степень
два, а $\deg_yQ_1\ge2$, поскольку $p_N\equiv2-y^2\pmod x$.

По модулю $x^{N+1}$ линейный коэффициент $Q_1$ обратим, а все
коэффициенты степеней $y^j$, $j\ge2$, нильпотентны.
Критерий Рассела--Сатайе показывает, что $G_N$ является
$R$-координатой. Но
\[
 [y^2]Q_1=-x\pmod{x^2},\qquad
 x^{N+1}\nmid[y^2]Q_1\quad(N\ge1).
\]
По критерию длины один $G_N$ дикая над $R$.
Лемма~\ref{lem:inner-wild} исключает ручность
$H_N-Q_x(c_N)$, а вычитание полинома от $x$ ручность не меняет.
Теорема Умирбаева--Ю даёт дикость $H_N$ в трёх переменных.
Поскольку $H_N$ входит в $\Phi_N$, последняя также дикая.
\end{proof}

Это доказательство относится ко всем $N$, а не выводится из проверки
конечного числа индексов. При $N=1$ его внутренний полином равен
$G_1=x^2z+by-xy^2$. Компьютерные проверки нормировки и коэффициентов
служат проверкой выкладок, но не заменяют лемму~\ref{lem:inner-wild}.

\section{Стандартная кремонова подгруппа}\label{sec:standard-subgroup}

Зафиксируем
\[
 \mathcal C_3(K)=\langle\operatorname{PGL}_4(K),\sigma_3\rangle,
 \qquad
 \sigma_3[X_0:X_1:X_2:X_3]
 =[X_1X_2X_3:X_0X_2X_3:X_0X_1X_3:X_0X_1X_2].
\]
В этом разделе произведения карт действуют справа налево.
По Блану--Хедену, $\Tame_3(K)\subset\mathcal C_3(K)$.
Их мономиальный критерий в размерности три утверждает, что
мономиальная бирациональная карта с матрицей показателей $M$ входит
в $\mathcal C_3(K)$ тогда и только тогда, когда суммы всех столбцов
$M$ нечётны \cite[теорема~1.1 и предложение~6.5]{BH2015}.

\begin{proposition}\label{prop:cremona-membership}
Все построенные выпрямители $\Phi_T,\Phi_H,\Phi_N$ принадлежат
$\mathcal C_3(K)$.
\end{proposition}
\begin{proof}
Рассмотрим подгруппу $\mathcal L\subset\Cr_3(K)$,
фиксирующую $x$ и порождённую картами
\[
 \begin{aligned}
 E_y(q)&=(x,y+q(x,z),z),& E_z(q)&=(x,y,z+q(x,y)),\\
 S&=(x,z,y),&B_r&=(x,x^ry,x^{-r}z),
 \end{aligned}
\]
где $q\in K[x,x^{-1},T]$, $r\in\Z$, и постоянными
диагональными картами $(x,cy,dz)$, $c,d\in K^\times$.

Все генераторы входят в $\mathcal C_3(K)$.
Для $B_r$ суммы столбцов матрицы показателей равны $(1,1,1)$.
Для $D=(x,x^{2s}y,z)$ они равны $(1+2s,1,1)$.
Выбрав $s\ge0$ так, чтобы $x^{2s}q$ был полиномом, получим
\[
 E_y(q)=D^{-1}E_y(x^{2s}q)D.
\]
Средняя карта ручная. Карты $E_z(q)$ получаются сопряжением с $S$;
остальные постоянные генераторы линейны.

В переменных $(u,v)$, при неизменном $x$, базовая карта раскладывается как
\[
 (u,v)\longmapsto\left(P_x(u)+mv,
 \frac{u-F_x(P_x(u)+mv)}m\right)
 =T_FD_mST_{P/m},
\]
где
\[
 \begin{aligned}
 T_{P/m}(u,v)&=(u,v+P_x(u)/m),&D_m(u,v)&=(mu,v/m),\\
 T_F(u,v)&=(u,v-F_x(u)/m),&S(u,v)&=(v,u).
 \end{aligned}
\]
При $m=bx^5$ каждый множитель принадлежит $\mathcal L$.
Начальная замена $(y,z)\mapsto(axy+z,-y)$ также является произведением
генераторов $\mathcal L$. Следовательно, $A_0\in\mathcal L$.

Для переноса положим $D_N=(x,y,x^Nz)$ и используем исправляющий
полином $\varphi_N(x,T)$ из конструкции полного компаньона. Тогда
\[
 \beta_N=SE_z(p_N)D_N,\qquad
 A_N=D_N^{-1}\bigl(E_z(-\varphi_N)A_0SE_z(p_N)\bigr)D_N.
\]
Карта $D_N$ нормализует $\mathcal L$, поскольку
\[
 \begin{aligned}
 D_N^{-1}E_y(q)D_N&=E_y(q(x,x^NT)),\\
 D_N^{-1}E_z(q)D_N&=E_z(x^{-N}q),\\
 D_N^{-1}SD_N&=B_NS,
 \end{aligned}
\]
а с $B_r$ и постоянными диагональными картами она коммутирует.
Те же формулы с $-N$ дают обратное включение.
Поэтому $A_N\in\mathcal L\subset\mathcal C_3(K)$.
Внешние карты, завершающие выпрямления, ручные; утверждение следует.
\end{proof}

\begin{remark}
При нечётном $N$ сама карта $D_N$ не проходит мономиальный критерий:
первая сумма столбца равна $1+N$. В доказательстве используется
\emph{нормализация} $\mathcal L$ картой $D_N$, а не принадлежность
$D_N$ группе $\mathcal C_3(K)$.
\end{remark}

\section{Стабильная ручность и границы вывода}\label{sec:stable-tameness}

\begin{proposition}\label{prop:stable-tameness}
Для каждого из построенных выпрямителей
\[
 \Phi_\ast\times\id_{\A^3}\in\Tame_6(K),
 \qquad \ast=T,H,N\ (N\ge1).
\]
\end{proposition}
\begin{proof}
Берсон--ван ден Эссен--Райт доказали, что подгруппа
$\Aut_K(\A^3)$, порождённая автоморфизмами, фиксирующими одну
координату, после добавления трёх переменных содержится в
$\Tame_6(K)$ в характеристике нуль
\cite[следствие~4.9]{BEW2012}.
Каждый $A_0,A_N$ фиксирует $x$.
Внешние множители наших выпрямителей ручные. Элементарные
автоморфизмы фиксируют две координаты, а аффинная группа порождена
элементарными линейными преобразованиями, масштабированиями одной
координаты и переносами одной координаты; каждый такой генератор фиксирует хотя бы одну
координату. Поэтому все $\Phi_\ast$ входят в указанную подгруппу.
\end{proof}

Оценка в три дополнительные переменные является достаточной;
минимальное число здесь не устанавливается. Объединяя доказательства,
получаем рабочий итог
\[
 \Phi_\ast\in
 \bigl(\mathcal C_3(K)\cap\Aut_K(\A^3)\bigr)
 \setminus\Tame_3(K),\qquad
 \Phi_\ast\times\id_{\A^3}\in\Tame_6(K).
\]
Это свойства \emph{выбранных карт}. Существование другого ручного
выпрямителя той же кривой ими не исключается.

Существование диких автоморфизмов в трёх переменных и диких координат
уже установлено в литературе. Принадлежность дикого автоморфизма Нагаты
группе $\mathcal C_3$ также доказана ранее
\cite[предложение~6.8]{BH2015}. Поэтому эти результаты не заявляются
как решение общей проблемы ручности, общей гипотезы стабильной ручности
или задачи о порождающих всей группы $\Cr_3$.
Содержание данного применения состоит в одновременной характеристике
явных выпрямителей конкретных кривых и в доказательстве для всех
индексов построенного семейства.
